\subsection{分裂半单李代数的自同构}

设 E 是一个交换群，\(A = \bigoplus_{\gamma \in E} A^{\gamma}\) 是一个 E-分次代数。对任何同态 \(\varphi\) （从群 E 到乘法群 \(k^{*}\) 的），用 \(f(\varphi)\) 表示从 A 到 A 的 k-线性映射，它在每个 \(A^{\gamma}\) 上的限制是比为 \(\varphi(\gamma)\) 的位似；显然 \(f(\varphi)\) 是分次代数 A 的自同构，且 f 是从群 \(\operatorname{Hom}(E, k^{*})\) 到分次代数 A 的自同构群的同态。

设 h 是 g 的一个分裂嘉当子代数，R 是 \((\mathfrak{g}, \mathfrak{h})\) 的根系。回顾 P(R) （相应地 Q(R) ）表示 R 的权格（相应地根权格）。置

\[
\mathrm{T} _ {\mathrm{P}} = \operatorname{Hom} (\mathrm{P} (\mathrm{R}), k ^ {*}) \quad \mathrm{T} _ {\mathrm{Q}} = \operatorname{Hom} (\mathrm{Q} (\mathrm{R}), k ^ {*}).
\]

我们可以把 \(g = g^{0} + \sum_{\alpha \in R} g^{\alpha}\) 视为一个 \(Q(R)\) -分次代数。前面的注记定义了一个从 \(T_{Q}\) 到 \(\operatorname{Aut}(\mathfrak{g}, \mathfrak{h})\) 的典则同态，它在本节中将记作 f。另一方面，从 \(Q(R)\) 到 \(P(R)\) 的典则单射定义了一个从 \(T_{P}\) 到 \(T_{Q}\) 的同态，将记作 q：

\[
\mathrm{T} _ {\mathrm{P}} \quad \stackrel {{q}} {{\longrightarrow}} \quad \mathrm{T} _ {\mathrm{Q}} \quad \stackrel {{f}} {{\longrightarrow}} \quad \operatorname{Aut} (\mathfrak {g}, \mathfrak {h}).
\]

若 \(s \in \operatorname{Aut}(\mathfrak{g}, \mathfrak{h})\) ，用 \(s^*\) 表示 \({}^t(s|_{\mathfrak{h}})^{-1}\) 在 \(\mathrm{Q}(\mathrm{R})\) 上的限制。则对一切 \(\varphi \in \mathrm{T}_{\mathrm{Q}}\) ，

\[
f (\varphi \circ s ^ {*}) = s ^ {- 1} \circ f (\varphi) \circ s.\tag{1}
\]

事实上，设 \(\gamma \in \mathrm{Q}(\mathrm{R})\) 与 \(x\in \mathfrak{g}^{\gamma}\) ；则 \(sx\in \mathfrak{g}^{s^{*}\gamma}\) ，且

\[
f (\varphi \circ s ^ {*}) x = (\varphi \circ s ^ {*}) (\gamma). x = s ^ {- 1} (\varphi (s ^ {*} \gamma) s x) = (s ^ {- 1} \circ f (\varphi) \circ s) (x).
\]

命题 2. 同态序列

\[
1 \longrightarrow \mathrm {T_ {Q}} \stackrel {f} {\longrightarrow} \mathrm{Aut} (\mathfrak {g}, \mathfrak {h}) \stackrel {\varepsilon} {\longrightarrow} \mathrm{A(R)} \longrightarrow 1
\]

是正合的。

\begin{enumerate}
\def\labelenumi{\alph{enumi})}
\item
  设 \(\varphi \in \operatorname{Ker} f\) 。则 \(\varphi(\alpha) = 1\) 对一切 \(\alpha \in \mathrm{R}\) 成立。由于 \(\mathrm{R}\) 生成群 \(\mathrm{Q}(\mathrm{R})\) ，\(\varphi\) 是 \(\mathrm{T}_{\mathrm{Q}}\) 的单位元。
\item
  设 \(\varphi \in \mathrm{T}_{\mathrm{Q}}\) 。\(f(\varphi)\) 在 \(\mathfrak{h} = \mathfrak{g}^0\) 上的限制是恒等映射，故
\end{enumerate}

\(\operatorname{Im} f \subset \operatorname{Ker} \varepsilon\) 。

\begin{enumerate}
\def\labelenumi{\alph{enumi})}
\setcounter{enumi}{2}
\tightlist
\item
  设 \(s \in \operatorname{Ker} \varepsilon\) 。则 \(s|_{\mathfrak{h}} = \mathrm{Id}_{\mathfrak{h}}\) 。对一切 \(\alpha \in R\) ，我们有 \(s(\mathfrak{g}^{\alpha}) = \mathfrak{g}^{\alpha}\) ，且存在 \(t_{\alpha} \in k^{*}\) 使 \(sx = t_{\alpha}x\) 对一切 \(x \in g^{\alpha}\) 成立。写下 \(s \in \operatorname{Aut}(\mathfrak{g})\) 这一条件，我们得到诸关系式
\end{enumerate}

\[
t _ {\alpha} t _ {- \alpha} = 1 \qquad \text{对一切 } \alpha \in \mathrm{R}
\]

\[
t _ {\alpha} t _ {\beta} = t _ {\alpha + \beta} \quad \text { when } \alpha , \beta , \alpha + \beta \in \mathrm{R}.
\]

在这些条件下，存在 \(\varphi \in \mathrm{T}_{\mathrm{Q}}\) 使 \(\varphi (\alpha) = t_{\alpha}\) 对一切 \(\alpha \in \mathrm{R}\) 成立（第六章 §1 第 6 节 命题 19 的推论 2）。则 \(s = f(\varphi)\) 。因此，\(\operatorname{Ker} \varepsilon \subset \operatorname{Im} f\) 。

\begin{enumerate}
\def\labelenumi{\alph{enumi})}
\setcounter{enumi}{3}
\tightlist
\item
  \(\operatorname{Aut}(\mathfrak{g},\mathfrak{h})\) 在 \(\varepsilon\) 下的像由 §2 第 2 节 定理 2 的推论知包含 \(\mathrm{W}(\mathrm{R})\) ，并由命题 1 知包含 \(\operatorname{Aut}(\mathrm{R},\mathrm{B})\) 。因此这个像等于 \(\mathrm{A}(\mathrm{R})\) 。
\end{enumerate}

推论 1. 设 \((\mathrm{B}, (X_{\alpha})_{\alpha \in \mathrm{B}})\) 是 \((\mathfrak{g}, \mathfrak{h})\) 的一个标架。设 \(\mathrm{G}\) 是使该标架不变的 \(s \in \operatorname{Aut}(\mathfrak{g}, \mathfrak{h})\) 组成的集合。则 \(\operatorname{Aut}(\mathfrak{g}, \mathfrak{h})\) 是 \(\mathrm{G}\) 与 \(\varepsilon^{-1}(\mathrm{W}(\mathrm{R}))\) 的半直积。

事实上，由命题 1，\(\mathrm{G} \cap \varepsilon^{-1}(\mathrm{W}(\mathrm{R})) = \{1\}\) ，且

\[
\operatorname{Aut} (\mathfrak {g}, \mathfrak {h}) = \mathrm{G}. \varepsilon^ {- 1} (\mathrm{W(R)})
\]

因为 \(\varepsilon\) 是满射（命题 2）。

推论 2. 群 \(\varepsilon^{-1}(\mathrm{W}(\mathrm{R}))\) 单可递地作用在 \((\mathfrak{g},\mathfrak{h})\) 的标架组成的集合上。

事实上，由 §4 第 4 节 定理 2，\(\operatorname{Aut}(\mathfrak{g},\mathfrak{h})\) 可递地作用在 \((\mathfrak{g},\mathfrak{h})\) 的标架组成的集合上。推论 2 于是由推论 1 得出。

推论 3. 设 B 是 R 的一个基。群 \(\operatorname{Ker} \varepsilon = f(\mathrm{T}_{\mathrm{Q}})\) 单可递地作用在 \((\mathfrak{g}, \mathfrak{h})\) 的形如 \((\mathrm{B}, (X_{\alpha})_{\alpha \in \mathrm{B}})\) 的标架组成的集合上。

这立即由命题 2 得出。

设 \(\alpha \in \mathrm{R}\) ，\(X_{\alpha} \in \mathfrak{g}^{\alpha}\) ，\(X_{-\alpha} \in \mathfrak{g}^{-\alpha}\) 满足 \([X_{\alpha}, X_{-\alpha}] = -H_{\alpha}\) 。我们已看到（§2 第 2 节 定理 2），对一切 \(t \in k^{*}\) ，初等自同构

\[
\theta_ {\alpha} (t) = e ^ {\mathrm{ad} t X _ {\alpha}} e ^ {\mathrm{ad} t ^ {- 1} X _ {- \alpha}} e ^ {\mathrm{ad} t X _ {\alpha}}
\]

在 \(\mathfrak{h}\) 上的限制是 \(s_{\alpha}\) 的转置；于是 \(\varepsilon(\theta_{\alpha}(t)) = s_{\alpha}\) ，从而 \(\theta_{\alpha}(t)\theta_{\alpha}(-1) \in \operatorname{Ker} \varepsilon\) 。

引理 1. 设 \(\alpha \in \mathrm{R}\) 与 \(t\in k^{*}\) 。设 \(\varphi\) 是同态 \(\lambda \mapsto t^{\lambda (H_{\alpha})}\) （从 \(\mathrm{Q(R)}\) 到 \(k^*\) 的）。则 \(f(\varphi) = \theta_{\alpha}(t)\theta_{\alpha}(-1)\) 。

设 \(\rho\) 是 \(\mathfrak{sl}(2,k)\) 在 \(\mathfrak{g}\) 上的与 \(X_{\alpha}\) 相伴的表示。设 \(\pi\) 是 \(\mathbf{SL}(2,k)\) 的与 \(\rho\) 相容的表示。采用 §1 第 5 节 的记号 \(\theta(t), h(t)\) 。由于 \(\rho(H) = \operatorname{ad} H_{\alpha}\) ，\(\mathfrak{g}^{\lambda}\) 的元素的权是 \(\lambda(H_{\alpha})\) （对 \(\rho\) 而言）。由 §2 第 2 节，\(\theta_{\alpha}(t)\theta_{\alpha}(-1) = \pi(\theta(t)\theta(-1)) = \pi(h(t))\) 。因此 \(\theta_{\alpha}(t)\theta_{\alpha}(-1)\) 在 \(\mathfrak{g}^{\lambda}\) 上的限制是比为 \(t^{\lambda(H_{\alpha})}\) 的位似（§1 第 5 节 命题 6），由此得引理。

命题 3. 复合同态

\[
\mathrm{T} _ {\mathrm{P}} \quad \xrightarrow {q} \quad \mathrm{T} _ {\mathrm{Q}} \quad \xrightarrow {f} \quad \operatorname{Aut} (\mathfrak {g}, \mathfrak {h})
\]

的像包含于 \(\mathrm{Aut}_e(\mathfrak{g})\) 。

设 B 是 R 的一个基。则 \((H_{\alpha})_{\alpha\in\mathrm{B}}\) 是 \(R^{\vee}\) 的一个基，且 \((H_{\alpha})_{\alpha\in\mathrm{B}}\) 在 \(h^{*}\) 中的对偶基是群 P(R) 的一个基。因此群 \(T_{P}\) 由诸同态 \(\lambda\mapsto t^{\lambda(H_{\alpha})}\;(t\in k^{*},\alpha\in\mathrm{B})\) 生成。若 \(\varphi\) 是这样一个同态到 Q(R) 上的限制，引理 1 证明 \(f(\varphi)\in\operatorname{Aut}_{e}(\mathfrak{g})\) ，由此得命题。

设 \(\bar{k}\) 是 \(k\) 的一个代数闭包。把每个自同构 \(s\) （\(\mathfrak{g}\) 的）对应到自同构 \(s \otimes 1\) （\(\mathfrak{g} \otimes_k \bar{k}\) 的）的映射是从 \(\mathrm{Aut}(\mathfrak{g})\) 到 \(\mathrm{Aut}(\mathfrak{g} \otimes_k \bar{k})\) 的单同态。我们用 \(\mathrm{Aut}_0(\mathfrak{g})\) 表示 \(\mathrm{Aut}(\mathfrak{g})\) 的正规子群，它是 \(\mathrm{Aut}_e(\mathfrak{g} \otimes_k \bar{k})\) 在这个同态下的逆像；这就是 \(\mathfrak{g}\) 的那样一些自同构组成的集合：它们在把基域从 \(k\) 扩张到 \(\bar{k}\) 后成为初等的。显然 \(\mathrm{Aut}_e(\mathfrak{g})\) 与 \(\bar{k}\) 的选取无关，且 \(\mathrm{Aut}_e(\mathfrak{g}) \subset \mathrm{Aut}_0(\mathfrak{g})\) 。群 \(\mathrm{Aut}_0(\mathfrak{g})\) 与 \(\mathrm{Aut}_e(\mathfrak{g})\) 可以不同（第七章 §13 第 1 节）。若 \(\mathfrak{h}\) 是 \(\mathfrak{g}\) 的一个嘉当子代数，置

\[
\operatorname{Aut} _ {e} (\mathfrak {g}, \mathfrak {h}) = \operatorname{Aut} _ {e} (\mathfrak {g}) \cap \operatorname{Aut} (\mathfrak {g}, \mathfrak {h}), \quad \operatorname{Aut} _ {0} (\mathfrak {g}, \mathfrak {h}) = \operatorname{Aut} _ {0} (\mathfrak {g}) \cap \operatorname{Aut} (\mathfrak {g}, \mathfrak {h}).
\]

引理 2. 设 \(\mathfrak{h}\) 是 \(\mathfrak{g}\) 的一个分裂嘉当子代数，且 \(s \in \operatorname{Aut}_0(\mathfrak{g}, \mathfrak{h})\) 。假设 \(s\) 在 \(\sum_{\alpha \in \mathrm{R}} \mathfrak{g}^\alpha\) 上的限制不以 1 为特征值。则 \(\varepsilon(s) = 1\) 。

通过扩张 \(k\) ，我们化归到 \(s \in \mathrm{Aut}_e(\mathfrak{g}, \mathfrak{h})\) 的情形。\(s - 1\) 的幂零空间的维数至少是 \(\dim \mathfrak{h}\) （第七章 §4 第 4 节 命题 9）。因此 \((s - 1)|_{\mathfrak{h}}\) 是幂零的。由于 \(s|_{\mathfrak{h}} \in \mathrm{A}(\mathrm{R}^\vee)\) ，\(s|_{\mathfrak{h}}\) 是有限阶的，从而是半单的（第五章附录命题 2）。于是 \((s - 1)|_{\mathfrak{h}} = 0\) ，这就证明了 \(\varepsilon(s) = 1\) 。

引理 3. (i) 设 \(m = (\mathrm{P}(\mathrm{R}) : \mathrm{Q}(\mathrm{R}))\) 。若 \(\varphi\) 是某个元素的 \(m\) 次幂（\(\mathrm{T}_{\mathrm{Q}}\) 中的元素），则 \(\varphi \in q(\mathrm{T}_{\mathrm{P}})\) 。

\begin{enumerate}
\def\labelenumi{(\roman{enumi})}
\setcounter{enumi}{1}
\tightlist
\item
  若 \(k\) 是代数闭的，则 \(q(\mathrm{T}_{\mathrm{P}}) = \mathrm{T}_{\mathrm{Q}}\) 。
\end{enumerate}

存在 P(R) 的一个基 \((\lambda_{1},\ldots,\lambda_{l})\) 以及整数 \(n_{1}\geq1,\ldots,n_{l}\geq1\) ，使 \((n_{1}\lambda_{1},\ldots,n_{l}\lambda_{l})\) 是 Q(R) 的一个基。我们有 \(m=n_{1}\ldots n_{l}\) 。设 \(\psi\in T_{Q}\) ，并置 \(\psi(n_{1}\lambda_{1})=t_{1},\ldots,\psi(n_{l}\lambda_{l})=t_{l}\) 。对 \(i=1,\ldots,l\) ，置 \(m_{i}=\) \(\prod_{j \neq i} n_j\) 。设 \(\chi\) 是 \(T_P\) 的元素，使 \(\chi(\lambda_1) = t_1^{m_1}, \ldots, \chi(\lambda_l) = t_l^{m_l}\) 。则

\[
\chi (n _ {i} \lambda_ {i}) = t _ {i} ^ {m _ {i} n _ {i}} = t _ {i} ^ {m} = (\psi^ {m}) (n _ {i} \lambda_ {i})
\]

于是 \(\chi|_{\mathrm{Q}(\mathrm{R})}=\psi^{m}\) 。这就证明了 (i)。若 k 是代数闭的，\(k^{*}\) 的每个元素都是 \(k^{*}\) 的某个元素的 m 次幂，故 \(T_{Q}\) 的每个元素都是 \(T_{Q}\) 的某个元素的 m 次幂；于是 (ii) 由 (i) 得出。

命题 4. 我们有 \(f(\mathrm{T}_{\mathrm{Q}}) \subset \mathrm{Aut}_0(\mathfrak{g}, \mathfrak{h})\) 与 \(\varepsilon^{-1}(\mathrm{W}(\mathrm{R})) = \mathrm{Aut}_0(\mathfrak{g}, \mathfrak{h})\) 。

\begin{enumerate}
\def\labelenumi{\alph{enumi})}
\tightlist
\item
  设 \(\varphi \in \mathrm{T}_{\mathrm{Q}}\) ，并设 \(\bar{k}\) 是 \(k\) 的一个代数闭包。由引理 3，\(\varphi\) 可以扩张为 \(\operatorname{Hom}(\mathrm{P}(\mathrm{R}), k^*)\) 的一个元素。由命题 3，
\end{enumerate}

\[
f (\varphi) \otimes 1 \in \operatorname{Aut} _ {e} (\mathfrak {g} \otimes_ {k} \bar {k}, \mathfrak {h} \otimes_ {k} \bar {k}).
\]

因此 \(f(\varphi) \in \mathrm{Aut}_0(\mathfrak{g},\mathfrak{h})\) ，且 \(\operatorname{Ker} \varepsilon \subset \mathrm{Aut}_0(\mathfrak{g},\mathfrak{h})\) 。

\begin{enumerate}
\def\labelenumi{\alph{enumi})}
\setcounter{enumi}{1}
\item
  \(\mathrm{Aut}_e(\mathfrak{g},\mathfrak{h})\) 在 \(\varepsilon\) 下的像包含 \(\mathrm{W(R)}\) （§2 第 2 节 定理 2 的推论）。鉴于 a)，我们看到 \(\varepsilon^{-1}(\mathrm{W(R)})\subset \mathrm{Aut}_0(\mathfrak{g},\mathfrak{h})\) 。
\item
  剩下要证 \(\mathrm{Aut}_0(\mathfrak{g},\mathfrak{h})\subset \varepsilon^{-1}(\mathrm{W}(\mathrm{R}))\) 。鉴于 b)，只需证明 \(\varepsilon (\mathrm{Aut}_0(\mathfrak{g},\mathfrak{h}))\cap \mathrm{Aut}(\mathrm{R},\mathrm{B})\) （其中 B 表示 R 的一个基）约化为 \{1\}。
\end{enumerate}

设 \(s \in \mathrm{Aut}_0(\mathfrak{g}, \mathfrak{h})\) 满足 \(\varepsilon(s) \in \mathrm{Aut}(\mathrm{R}, \mathrm{B})\) 。\(\mathrm{A}(\mathrm{R})\) 的由 \(\varepsilon(s)\) 生成的子群在 \(\mathrm{R}\) 上只有有限个轨道。设 \(\mathrm{U}\) 是这样一个轨道，其基数是 \(r\) ，并设 \(\mathfrak{g}^{\mathrm{U}} = \sum_{\beta \in \mathrm{U}} \mathfrak{g}^{\beta}\) 。设 \(\beta_1 \in \mathrm{U}\) ，并置 \(\beta_i = \varepsilon(s)^{i-1}\beta_1\) 对 \(1 \leq i \leq r\) ，于是 \(\mathrm{U} = \{\beta_1, \ldots, \beta_r\}\) 。设 \(X_{\beta_1}\) 是 \(\mathfrak{g}^{\beta_1}\) 的一个非零元素，并置 \(X_{\beta_i} = s^{i-1}X_{\beta_1}\) 对 \(1 \leq i \leq r\) 。存在 \(c_{\mathrm{U}} \in k^*\) 使 \(s^r X_{\beta_1} = c_{\mathrm{U}}X_{\beta_1}\) ，于是 \(s^r X_{\beta_i} = c_{\mathrm{U}}X_{\beta_i}\) 对一切 \(i\) 成立，从而 \(s^r |_{\mathfrak{g}^{\mathrm{U}}} = c_{\mathrm{U}}.1\) 。设 \(\varphi \in T_{\mathrm{Q}}\) ，并设 \(s' = s \circ f(\varphi)\) ，由 \(a)\) 它是 \(\mathrm{Aut}_0(\mathfrak{g}, \mathfrak{h})\) 的元素。我们有 \(s'^r |_{\mathfrak{g}^{\mathrm{U}}} = c_{\mathrm{U}}'.1\) ，其中

\[
c _ {\mathrm{U}} ^ {\prime} = c _ {\mathrm{U}} \prod_ {i = 1} ^ {r} \varphi (\beta_ {i}) = c _ {\mathrm{U}} \varphi \left(\sum_ {i = 1} ^ {r} \beta_ {i}\right).
\]

置 \(\mathrm{B} = \{\alpha_1, \ldots, \alpha_l\}\) 及 \(\sum_{i=1}^{r} \beta_i = \sum_{j=1}^{l} m_j^{\mathrm{U}} \alpha_j\) 。由于 \(\varepsilon(s) \in \mathrm{Aut}(\mathrm{R}, \mathrm{B})\) ，诸 \(m_j^{\mathrm{U}}\) 是同号且不全为零的整数。我们有

\[
c _ {\mathrm{U}} ^ {\prime} = c _ {\mathrm{U}} \prod_ {j = 1} ^ {l} \varphi (\alpha_ {j}) ^ {m _ {j} ^ {\mathrm{U}}}.
\]

现在可以选取 \(\varphi\) ，使对每个轨道 U 都有 \(c_{\mathrm{U}}^{\prime} \neq 1\) ；事实上，这归结为选取元素 \(\varphi(\alpha_1) = t_1, \ldots, \varphi(\alpha_l) = t_l\) （\(k^*\) 中的），它们不被有限个 \(t_1, \ldots, t_l\) 的不恒为零的多项式所零化。对这样的 \(\varphi\) 的选取，由引理 2 有 \(\varepsilon(s') = 1\) ，故

\[
\varepsilon (s) = \varepsilon (s ^ {\prime}) \varepsilon (f (\varphi)) ^ {- 1} = 1.
\]

推论. 设 B 是 R 的一个基。群 \(\operatorname{Aut}(\mathfrak{g},\mathfrak{h})\) 同构于群 \(\operatorname{Aut}(\mathrm{R},\mathrm{B})\) 与 \(\operatorname{Aut}_{0}(\mathfrak{g},\mathfrak{h})\) 的半直积。

这由命题 1、命题 2 的推论 1 及命题 4 得出。

注. 设 \(\varepsilon^{\prime},\varepsilon^{\prime\prime}\) 是 \(\varepsilon\) 到 \(\operatorname{Aut}_{0}(\mathfrak{g},\mathfrak{h}),\operatorname{Aut}_{e}(\mathfrak{g},\mathfrak{h})\) 上的限制。设 \(f^{\prime}\) 是从 \(T_{P}\) 到 \(\operatorname{Aut}_{e}(\mathfrak{g},\mathfrak{h})\) 的、由 f 经由从 \(Q(R)\) 到 \(P(R)\) 的典则单射诱导的同态。在前面我们已建立了如下的交换图：

\[
\begin{array}{c c c c c c c c c} 1 & \longrightarrow & \mathrm{T} _ {\mathrm{Q}} & \stackrel {{f}} {{\longrightarrow}} & \mathrm{Aut} (\mathfrak {g}, \mathfrak {h}) & \stackrel {{\varepsilon}} {{\longrightarrow}} & \mathrm{A} (\mathrm{R}) & \longrightarrow & 1 \\ & & \uparrow & & \uparrow & & \uparrow & & \\ 1 & \longrightarrow & \mathrm{T} _ {\mathrm{Q}} & \stackrel {{f}} {{\longrightarrow}} & \mathrm{Aut} _ {0} (\mathfrak {g}, \mathfrak {h}) & \stackrel {{\varepsilon^ {\prime}}} {{\longrightarrow}} & \mathrm{W} (\mathrm{R}) & \longrightarrow & 1 \\ & & \uparrow_ {q} & & \uparrow & & \uparrow & & \\ & & \mathrm{T} _ {\mathrm{P}} & \stackrel {{f ^ {\prime}}} {{\longrightarrow}} & \mathrm{Aut} _ {e} (\mathfrak {g}, \mathfrak {h}) & \stackrel {{\varepsilon^ {\prime \prime}}} {{\longrightarrow}} & \mathrm{W} (\mathrm{R}) & \longrightarrow & 1 \end{array}
\]

其中除 q 外的竖直箭头表示典则单射。我们已看到（命题 2 与命题 4），前两行是正合的。在第三行中，同态 \(\varepsilon''\) 是满射（§2 第 2 节 定理 2 的推论）；可以证明其核是 \(f'(\mathrm{T}_{\mathrm{P}})\) （§7 习题 26 d)）。

